\documentclass{article}
\input{../../../../lib.sty}

\usepackage[backend=biber, style=authoryear-comp, sorting=nyt]{biblatex}
\addbibresource{ref.bib}

\title{\texttt{fn quantile\_cnd}}
\author{\AishwaryaRamasethuAuthor, \YuJuKuAuthor, \JordanAwanAuthor, \MichaelShoemateAuthor}
\begin{document}

\maketitle
\contrib

Compute the quantile of a canonical noise distribution, as specified by a tradeoff function \texttt{f}.

\section{Hoare Triple}

\subsection*{Precondition}
\subsubsection*{Compiler-verified}
\begin{itemize}
    \item Argument \texttt{uniform} of type \texttt{RBig}
    \item Argument \texttt{f}, a fallible function from \texttt{RBig} to \texttt{RBig}
    \item Argument \texttt{c} of type \texttt{RBig}
\end{itemize}

\subsubsection*{Caller-verified}
\begin{itemize}
    \item Argument \texttt{uniform} is in $[0, 1]$
    \item When successful, \texttt{f} evaluates a fixed symmetric nontrivial tradeoff function $f$ exactly.
    \item Argument \texttt{c} is the fixed point of $f$, with $0\leq c<1/2$.
\end{itemize}

\subsection*{Pseudocode}

\lstinputlisting[language=Python,firstline=2]{./pseudocode/quantile_cnd.py}

\subsection*{Postcondition}

\begin{theorem}
    \label{postcondition}
    If \texttt{uniform} $\in (0,1)$, the pseudocode evaluates the quantile function $F^{-1}_f(u)$
    as defined in Proposition~F.6 of \cite{awan2023canonical}, or propagates an error from evaluating $f$.
    If $\texttt{uniform} \in \{0,1\}$, the pseudocode returns \texttt{None}.
\end{theorem}

\begin{proof}

We start by defining $F_f(\cdot)$:
\begin{definition}[\protect{\cite[Definition~3.7]{awan2023canonical}}]
    \label{def1}  % Def 3.7
    Let $f$ be a symmetric nontrivial tradeoff function, and let {$c\in [0,1/2)$} be the unique fixed point of $f$: $f(c)=c$.
    We define $F_f:\mathbb{R}\rightarrow \mathbb{R}$ as
    \begin{equation}
        F_f(x) = \begin{cases}
            f(1-F_f(x+1))&x<-1/2\\
            c \cdot (1/2-x) + (1-c)(x+1/2)&-1/2\leq x\leq 1/2\\
            1-f(F_f(x-1))&x>1/2.\\
        \end{cases}
    \end{equation}
\end{definition}

The preconditions for \texttt{quantile\_cnd} satisfy the preconditions for this definition.
The quantile function $F^{-1}_f(u)$ is defined in the following lemma.
\begin{proposition}[\protect{\cite[Proposition~F.6]{awan2023canonical}}]
    \label{prop-f-6}
    The quantile function $F_f^{-1}:(0,1)\rightarrow \mathbb{R}$ for $F_f$ can be expressed as
    \begin{equation}
        F_f^{-1}(u) = \begin{cases}
            F_f^{-1}(1-f(u))-1&u<c\\
            \frac{u-1/2}{1-2c}&c\leq u\leq 1-c\\
            F_f^{-1}(f(1-u))+1&u>1-c,
        \end{cases}
    \end{equation}

    where $c$ is the unique fixed point of $f$.
    Furthermore, the recurrence terminates after finitely many steps for
    every $u\in(0,1)$.
\end{proposition}

Let $u_0$ denote the original input and $u$ the current value of
\texttt{uniform}. Endpoint inputs immediately return \texttt{None};
Proposition~\ref{prop-f-6} only specifies the quantile on the open interval.

For an interior input, the loop invariant is
\[
    F_f^{-1}(u_0)=\texttt{offset}+F_f^{-1}(u).
\]
It holds initially since \texttt{offset} is zero.
In the lower and upper tails, the first and third cases of
Proposition~\ref{prop-f-6} preserve the invariant by subtracting or
adding one to \texttt{offset}, respectively.
In the central interval the second case gives the exact returned value.
The denominator is positive since $c<1/2$.

For a nontrivial tradeoff, convexity and $f(v)\leq1-v$ imply
$f(v)<1-v$ for $0<v<1$: equality at an interior point would force
equality throughout the interval.
In a tail, the argument $v$ to $f$ lies in $(0,c)$, so
$f(v)\geq f(c)=c>0$ by nonincrease and $f(v)<1-v<1$.
Consequently the next uniform value is interior and differs from the
current value. If $c=0$ there are no interior tail cases.
Thus the endpoint, zero-denominator and stationary-value guards
cannot interrupt a valid interior recurrence.
Each iteration makes the same step as the proposition, so its finite
termination result applies. All arithmetic is rational and exact.
Errors evaluating $f$ propagate immediately.

\end{proof}

\printbibliography
\end{document}
